[V.] Art. $[7.]$ Hier bezeichnet $\mathrm{Dec.}\frac mM$, wie in Artikel $1$ des vorhergehenden Bruchstücks [IV], den Wert von
\[
\Sigma(-1)^X\Sigma(-1)^Y\mathrm{Int.}\,p
\]
wobei die Summation über alle ganzen Zahlen $X$ und $Y$ erstreckt wird, für die die zugehörigen $\xi$ und $\eta$ innerhalb der Grenzen $0$ und $\tfrac12$ liegen.

Nach der Angabe im Text wurden die Formeln für den Dezidenten in den Artikeln $7$ und $15$ auf zwei besonderen
\SourceLeaf{388}
Wegen gefunden. Um diese Erläuterungen jedoch nicht zu sehr auszudehnen, werden sie hier aus einer gemeinsamen Quelle abgeleitet.

Nimmt $X$ einen beliebigen bestimmten ganzzahligen Wert an, so sei $Y^{*}$ der kleinere und $Y^{**}$ der größere der beiden Werte von $Y$, die den Grenzwerten von $\xi$, $\eta$ entsprechen. Die zu $Y^{*}$ und $Y^{**}$ gehörenden Werte von $p$ seien $p^{*}$ und $p^{**}$. Ebenso wie $Y^{*}$ und $Y^{**}$ können sie nicht gleich sein, weil die Summe $\Sigma$ sich nicht über die Grenzwerte von $\xi$ und $\eta$ erstreckt.

Führt man wie in den beiden vorhergehenden Aufsätzen [III] und [IV] die Summation über alle beim selben $X$ auftretenden Werte von $Y$ aus, setzt dabei für die Anzahl der zwischen $Y'$ und $Y''$ liegenden ungeraden Zahlen $[\tfrac12Y''-\tfrac12]-[\tfrac12Y'-\tfrac12]$ und fügt die auf die Grenzen $\xi=0$ und $=\tfrac12$ bezogenen Intensoren zweimal mit entgegengesetzten Zeichen hinzu, so erhält man für $\Sigma(-1)^Y\mathrm{Int.}\,p$ den aus sieben Teilen bestehenden Ausdruck
\[
\begin{gathered}
-\Sigma[-\mathrm{Int.}(p-\mu\omega i)+\mathrm{Int.}(p+\mu\omega i)]\\
\text{wobei alle }p\text{ aufzunehmen sind, für die }[Y]\text{ gerade,}\\
x\text{ oder }y\text{ ganz, einschl. }\xi=0\text{ und }\tfrac12,\quad
\text{ausschl. }\eta=0\text{ und }\tfrac12
\end{gathered}
\]
\[
\left.\begin{aligned}
&-\mathrm{Int.}(p^{*}-\mu\omega i)\text{ wenn }[Y^{*}]\text{ gerade}\\
&+\mathrm{Int.}(p^{**}+\mu\omega i)\text{ wenn }[Y^{**}]\text{ gerade}
\end{aligned}\right\}\ \xi=0\text{ oder }\tfrac12,\quad0<\eta<\tfrac12
\]
\[
\left.\begin{aligned}
&-\mathrm{Int.}(p^{*}+\mu\omega i)\text{ wenn }[Y^{*}]\text{ gerade}\\
&+\mathrm{Int.}(p^{**}-\mu\omega i)\text{ wenn }[Y^{**}]\text{ gerade}
\end{aligned}\right\}\ 0<\xi<\tfrac12,\quad\eta=0\text{ oder }\tfrac12
\]
\[
\left.\begin{aligned}
&-\mathrm{Int.}(p^{*}+\mu\omega i)\text{ wenn }[Y^{*}]\text{ gerade}\\
&+\mathrm{Int.}(p^{**}-\mu\omega i)\text{ wenn }[Y^{**}-\omega]\text{ gerade}
\end{aligned}\right\}\ \xi=0\text{ oder }\tfrac12,\quad\eta=0\text{ oder }\tfrac12
\]
der, mit $(-1)^X$ multipliziert und über alle ganzzahligen $X$ summiert, den Dezidenten $\frac mM$ ergibt.

Aus dem ersten Teil des Ausdrucks entsteht so für die nach Artikel $2$, Vorschrift I gebildeten $\varepsilon$
\[
\begin{aligned}
&\Sigma\varepsilon\text{ wobei }X\text{ ganz, }[Y]\text{ gerade}\\
&-\Sigma\varepsilon\text{ wobei außerdem }y\text{ ganz, }[x]\text{ gerade}
\end{aligned}
\]
Für die folgenden Teile kann man
\[
\left.\begin{aligned}
&-(B)\mathrm{Int.}(p-B\mu\omega i)\text{ wenn }\xi=0\\
&+(B)\mathrm{Int.}(p+B\mu\omega i)\text{ wenn }\xi=\tfrac12
\end{aligned}\right\}\ X\text{ ganz, }[Y]\text{ gerade, }0<\eta<\tfrac12
\]
\[
\left.\begin{aligned}
&-(A)\mathrm{Int.}(p+A\mu\omega i)\text{ wenn }\eta=0\\
&+(A)\mathrm{Int.}(p-A\mu\omega i)\text{ wenn }\eta=\tfrac12
\end{aligned}\right\}\ X\text{ ganz, }[Y]\text{ gerade, }0<\xi<\tfrac12
\]
\[
\begin{gathered}
-\tfrac12[(A')+(B')]\mathrm{Int.}(p+A'\mu\omega i)\\
\text{wenn }X\text{ ganz, }[Y+A'\omega]\text{ gerade, }\xi\text{ und }\eta=0\text{ oder }\tfrac12
\end{gathered}
\]
setzen. Dabei ist $A'=+A$ oder $-A$, wenn $\eta=0$ oder $\tfrac12$, und $B'=+B$ oder $-B$, wenn $\xi=0$ oder $\tfrac12$; zum Beispiel verwendet man die Bezeichnung $(A):+1$ oder $-1$, je nachdem $A$ positiv oder negativ ist.

Multipliziert man mit $(-1)^X$, führt die Summation über $X$ aus und lässt dabei in diesen Ausdrücken, und zwar im
\[
\begin{aligned}
\text{ersten}\quad&p-B\mu\omega i,\ P-BD\omega i,\ \pi-AB\omega i-BB\omega,\ X,\ [Y]\\
&\text{jeweils in }ip,\ iP,\ i\pi,\ -Y,\ [X]\\[3pt]
\text{zweiten}\quad&p+B\mu\omega i,\ P+BD\omega i,\ \pi+AB\omega i+BB\omega,\ X,\ [Y]\\
&\dotstwo\quad p,\ P,\ \pi\quad X,\ [Y]\\[3pt]
\text{dritten}\quad&p+A\mu\omega i,\ P+AD\omega i,\ \pi+AA\omega i+AB\omega,\ X,\ [Y]\\
&\dotstwo\quad p,\ P,\ \pi\quad X,\ [Y]\\[3pt]
\text{vierten}\quad&p-A\mu\omega i,\ P-AD\omega i,\ \pi-AA\omega i-AB\omega,\ X,\ [Y]\\
&\dotstwo\text{ in }im-ip,\ iM-iP,\ i-i\pi,\ Y-B,\ A-1-[X]
\end{aligned}
\]

\SourceLeaf{389}
die angegebenen Größen ineinander übergehen und bezeichnet das Ergebnis aus dem fünften Ausdruck mit $Q_1$, so erhält man
\[
\begin{aligned}
&-\Sigma(-1)^Y(B)\mathrm{Int.}\,ip,\\
&\qquad\text{wobei }Y\text{ ganz, }[X]\text{ gerade, }\eta=\omega,\ 0<\xi<\tfrac12\\
&+\Sigma(-1)^X(B)\mathrm{Int.}\,p,\\
&\qquad\text{wobei }X\text{ ganz, }[Y]\text{ gerade, }\xi=\tfrac12+\omega,\ 0<\eta<\tfrac12\\
&-\Sigma(-1)^X(A)\mathrm{Int.}\,p,\\
&\qquad\text{wobei }X\text{ ganz, }[Y]\text{ gerade, }\eta=\omega,\ 0<\xi<\tfrac12\\
&+\Sigma(-1)^Y(A)\mathrm{Int.}(im-ip),\\
&\qquad\text{wobei }Y\text{ ganz, }[X]\text{ gerade, }\xi=\tfrac12+\omega,\ 0<\eta<\tfrac12\\
&+Q_1
\end{aligned}
\]

Die Untersuchung der einzelnen Fälle zeigt, dass unter der Voraussetzung $M\equiv1\ (\text{mod. }2+2i)$
\[
\begin{aligned}
Q_1={}&-\mathrm{Int.}\,\mu\omega i\text{ ist, wenn }M\text{ im }1.\text{ Quadranten liegt}\\
&-\mathrm{Int.}(\tfrac12mi+\mu\omega i),\quad\text{wenn }M\text{ im }2.\text{ Quadr. und }\frac{M-1}{2}\text{ gerade}\\
&+\mathrm{Int.}(\tfrac12mi-\mu\omega i),\quad\text{wenn }M\text{ im }4.\text{ Quadr. und }\frac{M-1}{2}\text{ gerade}
\end{aligned}
\]
wobei man eine komplexe Zahl gerade oder ungerade nennt, je nachdem sie durch $2$ teilbar ist oder nicht.

Bei Anwendung der in den Vorschriften II und IV bestimmten $\varepsilon$ ergibt sich somit
\[
\begin{aligned}
\mathrm{Dec.}\frac mM={}&\mathrm{I},& &\Sigma\varepsilon,\ \text{wobei }X\text{ ganz, }[Y]\text{ gerade}\\
&&&-4\Sigma\varepsilon,\ \text{wobei außerdem }y\text{ ganz, }[x]\text{ gerade}\\
&\mathrm{II},&&-\Sigma\varepsilon\mathrm{Int.}\,ip,\ \text{wobei }Y\text{ ganz, }[X]\text{ gerade}\\
&\mathrm{IV},&&+\Sigma\varepsilon\mathrm{Int.}\,p,\ \text{wobei }X\text{ ganz, }[Y]\text{ gerade}\\
&\mathrm{II},&&+\Sigma\varepsilon\mathrm{Int.}\,p,\ \text{wobei }X\text{ ganz, }[Y]\text{ gerade}\\
&\mathrm{IV},&&+\Sigma\varepsilon\mathrm{Int.}(im-ip),\ \text{wobei }Y\text{ ganz, }[X]\text{ gerade}\\
&&&+Q_1
\end{aligned}
\]

Eine andere Form des Ausdrucks für den Dezidenten erhält man, wenn man zuerst über $X$ summiert und die Anzahl der zwischen $X'$ und $X''$ liegenden ungeraden Zahlen durch $[\tfrac12X''+\tfrac12]-[\tfrac12X'+\tfrac12]$ darstellt, nämlich
\[
\begin{aligned}
\mathrm{Dec.}\frac mM={}&\mathrm{I},& &\Sigma\varepsilon,\ \text{wobei }Y\text{ ganz, }[X]\text{ ungerade}\\
&&&-4\Sigma\varepsilon,\ \text{wobei außerdem }y\text{ ganz, }[x]\text{ gerade}\\
&\mathrm{II},&&-\Sigma\varepsilon\mathrm{Int.}\,ip,\ \text{wobei }X\text{ ganz, }[Y]\text{ gerade}\\
&\mathrm{IV},&&+\Sigma\varepsilon\mathrm{Int.}\,p,\ \text{wobei }Y\text{ ganz, }[X]\text{ ungerade}\\
&\mathrm{II},&&+\Sigma\varepsilon\mathrm{Int.}\,p,\ \text{wobei }Y\text{ ganz, }[X]\text{ ungerade}\\
&\mathrm{IV},&&+\Sigma\varepsilon\mathrm{Int.}(im-ip),\ \text{wobei }X\text{ ganz, }[Y]\text{ ungerade}\\
&&&+Q_2
\end{aligned}
\]
\[
\begin{aligned}
Q_2={}&-\mathrm{Int.}(-\mu\omega),\quad\text{wenn }M\text{ im }2.\text{ Quadr.}\\
&+\mathrm{Int.}(\tfrac12m-\mu\omega),\quad\text{wenn }M\text{ im }1.\text{ Quadr. und }\frac{M-1}{2}\text{ ungerade}\\
&-\mathrm{Int.}(\tfrac12m+\mu\omega),\quad\text{wenn }M\text{ im }3.\text{ Quadr. und }\frac{M-1}{2}\text{ ungerade}
\end{aligned}
\]

\SourceLeaf{390}
Führt man die Summation über $Y$ zuerst aus, wählt aber die zweite soeben verwendete Methode, um die Anzahl der zwischen zwei Werten liegenden ungeraden Zahlen zu bestimmen, so ergibt sich
\[
\begin{aligned}
\mathrm{Dec.}\frac mM={}&\mathrm{I},& &\Sigma\varepsilon,\ \text{wobei }X\text{ ganz, }[Y]\text{ ungerade}\\
&&&-4\Sigma\varepsilon,\ \text{wobei außerdem }y\text{ ganz, }[x]\text{ gerade}\\
&\mathrm{II},&&-\Sigma\varepsilon\mathrm{Int.}\,ip,\ \text{wobei }Y\text{ ganz, }[X]\text{ ungerade}\\
&\mathrm{IV},&&+\Sigma\varepsilon\mathrm{Int.}\,p,\ \text{wobei }X\text{ ganz, }[Y]\text{ ungerade}\\
&\mathrm{II},&&+\Sigma\varepsilon\mathrm{Int.}\,p,\ \text{wobei }X\text{ ganz, }[Y]\text{ ungerade}\\
&\mathrm{IV},&&+\Sigma\varepsilon\mathrm{Int.}(im-ip),\ \text{wobei }Y\text{ ganz, }[X]\text{ ungerade}\\
&&&+Q_3
\end{aligned}
\]
\[
\begin{aligned}
Q_3={}&-\mathrm{Int.}(-\mu\omega i),\quad\text{wenn }M\text{ im }3.\text{ Quadr.}\\
&-\mathrm{Int.}(\tfrac12mi+\mu\omega i),\quad\text{wenn }M\text{ im }2.\text{ Quadr. und }\frac{M-1}{2}\text{ ungerade}\\
&+\mathrm{Int.}(\tfrac12mi-\mu\omega i),\quad\text{wenn }M\text{ im }4.\text{ Quadr. und }\frac{M-1}{2}\text{ ungerade}
\end{aligned}
\]

Summiert man auch zuerst über $X$ und verwendet die erste Darstellung der Anzahl der zwischen zwei Werten liegenden ungeraden Zahlen, so erhält man die in [V], Artikel $15$ angegebene Form für den Dezidenten. Dabei kann die Größe $q$ auch durch folgende Gleichung definiert werden:
\[
\begin{aligned}
q={}&-\mathrm{Int.}\,\mu\omega,\quad\text{wenn }M\text{ im }4.\text{ Quadr.}\\
&+\mathrm{Int.}(\tfrac12m-\mu\omega),\quad\text{wenn }M\text{ im }1.\text{ Quadr. und }\frac{M-1}{2}\text{ gerade}\\
&-\mathrm{Int.}(\tfrac12m+\mu\omega),\quad\text{wenn }M\text{ im }3.\text{ Quadr. und }\frac{M-1}{2}\text{ gerade}
\end{aligned}
\]

Die Vereinigung dieser vier Ausdrücke für den Dezidenten ergibt das in [V], Artikel $7$ aufgestellte Resultat. Denn $\mathrm{Int.}\,ip-\mathrm{Int.}\,p$ ist gleich $3$ für gerades $[x]$ und ungerades $[y]$, sonst aber gleich $1$; ferner ist $\mathrm{Int.}(im-ip)+\mathrm{Int.}\,p$ gleich $0$, wenn $[x]$ und $[y]$ gerade sind, in allen übrigen Fällen aber gleich $4$.

[V.] Art. $[7.]$ Die erste Tafel zum Beispiel gibt in der ersten Spalte die zu jedem ganzzahligen $P$ gehörenden Werte von $\frac{37\pdot P}{M}$ oder $37(\xi+\eta i)$, wenn $0<\xi<\tfrac12$, $0<\eta<\tfrac12$; in der zweiten $\frac{37\pdot Pm}{M}$ oder $37\pdot p$; in der dritten die größte in $p$ enthaltene ganze Zahl; und in der vierten $\pm\mathrm{Int.}\,p$, wobei das obere Zeichen gilt, wenn $P$ durch $1+i$ teilbar ist, das untere, wenn $P$ nicht durch $1+i$ teilbar ist.

[V.] Art. $[9.]$. $[12.]$ Die verbesserte Bezeichnung der $\theta$ ist nur bei der zweiten und dritten Klasse, Artikel $9$. $10$, angedeutet, aber auch auf die erste und vierte, Artikel $8$. $11$, auszudehnen. Demnach hat ein $\theta^\lambda=T+Ui$ denjenigen Index $\lambda,=0,1,2$ oder $3$, für den die durch die Gleichungen

\SourceLeaf{391}
\[
\begin{aligned}
i^\lambda M&=\mathfrak A+\mathfrak B i,\quad i^{-\lambda}\mu=\rho+\sigma i\\
\sigma\varphi^0&=-\text{Koeff. Imag. }\theta^\lambda(a-bi)=+bT-aU\\
\sigma\Phi^0&=+\text{Koeff. Imag. }\theta^\lambda(\mathfrak A-\mathfrak B i)=-\mathfrak B T+\mathfrak A U
\end{aligned}
\]
bestimmten Größen $\varphi^0$ und $\Phi^0$ zwischen $0$ und $\tfrac12$ liegen.

Um nach den Andeutungen in Artikel $9$ $(3)$ zu beweisen, dass, wenn $T$, $U$ zwei ganze reelle Zahlen sind, die die soeben aufgestellten Bedingungen erfüllen, $T+Ui$ auch in der bei einer der vier Kombinationen, Artikel $8$. $11$, bestimmten Gesamtheit $\theta^\lambda$ liegt, bezeichne man mit $\varphi'$, $\Phi'$ diejenigen komplexen ganzen Zahlen, für die die Gleichung
\[
T+Ui=\varphi'i^\lambda M+\Phi'm
\]
gilt und für die eine der vier Größen $\pm\frac{\varphi'-\varphi^0}{m}$, $\pm i\frac{\varphi'-\varphi^0}{m}$ so beschaffen ist, dass ihr reeller Teil und der Koeffizient ihres imaginären Teils zwischen $0$ und $\tfrac12$ liegen (\textit{Theoria residuorum biquadr. artt.} $45,46$). Die betreffende Größe ist dann, wie man aus der Untersuchung der sechzehn einzelnen Fälle in den vier Kombinationen leicht erkennt, $\frac pm$, und die ihr entsprechende Größe unter $\pm\frac{\Phi'-\Phi^0}{m}$, $\pm i\frac{\Phi'-\Phi^0}{M}$ ist $\frac PM$, weil $\frac{\Phi'-\Phi^0}{M}=-\frac{\varphi'-\varphi^0}{m}i^\lambda$.

Aus dieser Darstellung der Größen $\frac pm$ oder $\frac PM$ folgt auch, dass $\mathrm I,\ \Sigma\varepsilon$ von allen aus $\Theta+\Theta'+\Theta''+\Theta'''$ besteht. Dabei bezeichnet $\Theta^\lambda$ die Summe derjenigen $\varepsilon$, die für jeden Gitterpunkt $\theta$ innerhalb des Parallelogramms $0,\ \tfrac12m,\ \tfrac12m+\tfrac12i^\lambda M,\ \tfrac12i^\lambda M,$
\begin{itemize}
\item[] $=+1$ zu setzen sind, wenn $\theta$ durch $1+i$ teilbar und $\text{Koeff. Imag. }\mu i^{-\lambda}$ positiv ist oder wenn keine Bedingung gilt; dagegen
\item[] $=-1$, wenn nur eine gilt.
\end{itemize}

[V.] Art. $13$. Die Bestimmung von $\Theta^\lambda$ kann entweder mit den vier oben für $\mathrm{Dec.}\frac mM$ verwendeten Summationsarten oder, was im Wesentlichen dasselbe ist, mit den in [II], Artikel $11$ angedeuteten Methoden ausgeführt werden. Dabei sind die vier Konstruktionen zugrunde zu legen, die entstehen, wenn man die Punkte, deren $\theta$ ein Vielfaches von $1+i$ ist, jeweils mit den Punkten $\theta+1$, $\theta+i$, $\theta-1$ und $\theta-i$ verbindet.

Ordnet man auf der Begrenzung des zuvor genannten Parallelogramms allen Punkten ein $\theta$ zu, für die der reelle oder imaginäre Teil von $\theta$ eine ganze Zahl ist, bezeichnet mit $\theta^0$ die nächste durch $1+i$ teilbare ganze Zahl bei $\theta$ und mit $l$ die Ortsverschiebung von einem Punkt des geraden Begrenzungsstücks, das $\theta$ enthält, zu einem beliebigen nachfolgenden Punkt derselben Geraden, also zum Beispiel bei diesem Parallelogramm der Reihe nach die Größen $m,\ Mi^\lambda,\ -m,\ -Mi^\lambda$, und setzt
\[
\varepsilon=\pm1\text{ mit dem Zeichen des imaginären Teils von }\frac l{\theta-\theta^0}
\]
so ergibt die Vereinigung der auf die eine oder andere Weise erhaltenen vier Resultate $4\Theta^\lambda=-\Sigma\varepsilon$.

Die gesonderte Bestimmung der zu den Eckpunkten gehörenden $\theta$ und $\varepsilon$ wird vermieden, wenn man dieses Parallelogramm durch ein anderes ersetzt, dessen Begrenzungen denjenigen des ersten
\SourceLeaf{392}
unendlich nahe sind und das die beiden Punkte $0$ und $\tfrac12m+\tfrac12i^\lambda M$ nicht einschließt. Die Begrenzung eines solchen Parallelogramms erhält man, indem man sie auf die positiven Seiten der Linien
\[
\begin{gathered}
0\dotsfour\tfrac12m,\quad\tfrac12m+\tfrac12i^\lambda M\dotsfour\tfrac12m,\\
\tfrac12m+\tfrac12i^\lambda M\dotsfour\tfrac12i^\lambda M,\quad 0\dotsfour\tfrac12i^\lambda M
\end{gathered}
\]
legt. Ordnet man den vier so entstandenen Geraden der Reihe nach die unendlich kleinen positiven Größen $\omega_1,\omega_2,\omega_3,\omega_4$ zu, so kann man für die auf ihnen liegenden Punkte $\theta$
\[
\begin{aligned}
\theta&=p=m(\xi+\omega_1i),\\
-\theta i^\lambda+\tfrac12mi^\lambda+\tfrac12M&=P=M(\xi+\omega_2i),\\
-\theta+\tfrac12m+\tfrac12Mi^\lambda&=p=m(\xi+\omega_3i),\\
\theta i^{-\lambda}&=P=M(\xi+\omega_4i)\quad\text{wenn }\lambda\text{ gerade}
\end{aligned}
\]
\[
\begin{aligned}
\theta&=p=m(\xi+\omega_1i),\\
\theta i^{1-\lambda}-\tfrac12mi^{1-\lambda}+\tfrac12M&=P=M(\tfrac12+\omega_2+\eta i),\\
\theta i+\tfrac12m-\tfrac12Mi^{1+\lambda}&=p=m(\tfrac12+\omega_3+\eta i),\\
\theta i^{-\lambda}&=P=M(\xi+\omega_4i)\quad\text{wenn }\lambda\text{ ungerade}
\end{aligned}
\]
setzen. Dabei überschreiten $\xi$ und $\eta$ zur Schließung der Figur teilweise auch den Bereich der reellen Werte von $0$ bis $\tfrac12$ um unendlich kleine Größen.

Bezeichnen $G$, $g$, $H$, $h$ die Summen der jeweils nach den Vorschriften II, III, IV, V (in den Artikeln $3$ bis $6$) gebildeten $\varepsilon$, umfassen $G'$ oder $G_{\prime}$ und $g'$ oder $g_{\prime}$ diejenigen $\varepsilon$, die für die beim zweiten Parallelogramm gegebenenfalls auftretenden unendlich kleinen Werte von $\xi$ gelten, im Übrigen aber jeweils nach den Vorschriften II und III gebildet sind, beziehen sich ferner $G''$ oder $G_{\prime\prime}$ und $g''$ oder $g_{\prime\prime}$ ebenso auf dieselben Vorschriften, jedoch auf die unendlich kleinen Werte von $\tfrac12-\xi$, und schließlich $H',h_{\prime},H'',h_{\prime\prime}$ jeweils auf die Vorschriften IV, V, IV, V und die unendlich kleinen Werte von $\tfrac12-\eta,\tfrac12-\eta,\eta,\eta$, so gilt
\[
\begin{aligned}
4\Theta^\lambda={}&-(g+g'+g'')-i^\lambda(G+G'+G'')\\
&+i^\lambda(g+g_{\prime}+g_{\prime\prime})+(G+G_{\prime}+G_{\prime\prime})\quad\text{wenn }\lambda\text{ gerade}\\
4\Theta^\lambda={}&-(g+g'+g'')-i^{1-\lambda}(H+H'+H'')\\
&-i^{\lambda-1}(h+h_{\prime}+h_{\prime\prime})+(G+G_{\prime}+G_{\prime\prime})\quad\text{wenn }\lambda\text{ ungerade}
\end{aligned}
\]

Für den Eckpunkt $\theta$ des Parallelogramms, der dem Punkt $0$ am nächsten liegt, bezeichne $\xi_1$ den zugehörigen Wert des $\xi$ der ersten Seite und $\xi_4$ den zugehörigen Wert des $\xi$ der vierten Seite, sodass
\[
\theta=m(\xi_1+\omega_1i)=i^\lambda M(\xi_4+\omega_4i)
\]
Dann ergibt sich dasjenige $\xi$, dem auf der ersten Seite oder ihrer Verlängerung ein Punkt $p$ mit reellem Teil gleich $0$ entspricht, aus der Gleichung
\[
(\mathrm{Real.}\,p=0),\quad\xi-\xi_1=\sigma a\mathfrak A\omega_1-\sigma\omega_4
\]
wobei die positiven Faktoren der unendlich kleinen positiven Größen durch die Einheit ersetzt sind und $\sigma$, $\mathfrak A$ die durch
\[
\rho+\sigma i=i^{-\lambda}(\alpha+\beta i),\quad\mathfrak A+\mathfrak B i=i^\lambda(A+Bi)
\]
bestimmten reellen Größen bezeichnen. Dieser Punkt $p$ liegt auf der ersten Seite selbst, wenn $\xi-\xi_1$ positiv ist, also, wenn man $\omega_1$ unendlich klein gegenüber $\omega_4$ annimmt, wenn $\sigma$ negativ ist. Der dem Punkt $p$ nächste Punkt $p^0$, dessen darstellende Zahl durch $1+i$ geteilt wird, ist der Punkt $0$; daher hat
\SourceLeaf{393}
$\mathrm{Imag.}\frac m{p-p^0}$ oder $\mathrm{Imag.}\frac1{\xi+\omega_1i}$ das Minuszeichen. Für $\mathrm{Real.}\,p=0$ erhält man somit
\[
\varepsilon=-1\text{ wenn }(\sigma)=-1,\quad\varepsilon=0\text{ wenn }(\sigma)=+1,\quad\text{d. h. }\varepsilon=-\tfrac12+\tfrac12(\sigma)
\]
und ebenso für $\mathrm{Imag.}\,p=0$
\[
\begin{gathered}
\xi-\xi_1=\sigma b\mathfrak B\omega_1-\sigma\omega_4,\quad\mathrm{Imag.}\frac m{p-p^0}=\mathrm{Imag.}\frac1{\xi+\omega_1i},\\
\varepsilon=-\tfrac12+\tfrac12(\sigma)\quad\text{also }g'=-1+(\sigma)
\end{gathered}
\]

Bezüglich der vierten Seite gilt $P^0=0$, $\mathrm{Imag.}\frac M{P-P^0}=\mathrm{Imag.}\frac1{\xi+\omega_4i}$,
\[
\begin{aligned}
&\text{also für }\mathrm{Real.}(i^\lambda P)=0;&&\xi-\xi_4=-\sigma a\mathfrak A\omega_4+\sigma\omega_1,\\
&&&\varepsilon=-\tfrac12+\tfrac12(\sigma a\mathfrak A)\\
&\text{und für }\mathrm{Imag.}(i^\lambda P)=0;&&\xi-\xi_4=-\sigma b\mathfrak B\omega_4+\sigma\omega_1,\\
&&&\varepsilon=-\tfrac12+\tfrac12(\sigma b\mathfrak B)
\end{aligned}
\]
also $G_{\prime}=-1+\tfrac12(\sigma a\mathfrak A)+\tfrac12(\sigma b\mathfrak B)$ oder, weil $\rho=a\mathfrak A+b\mathfrak B$, $G_{\prime}=-1+\tfrac12(\rho\sigma)+\tfrac12(\rho\sigma ab\mathfrak A\mathfrak B)$.

Der Teil $R_1^\lambda$ von $4\Theta^\lambda$, der aus dem unendlich nahe beim Punkt $0$ liegenden Stück der Begrenzung entsteht, ist somit
\[
R_1^\lambda=+G_{\prime}-g'=-(\sigma)+\tfrac12(\rho\sigma)+\tfrac12(\rho\sigma ab\mathfrak A\mathfrak B)
\]

Durch ähnliche Überlegungen findet man für die Teile $R_2^\lambda$, $R_3^\lambda$, $R_4^\lambda$, die jeweils dieselbe Beziehung zu den Punkten $\tfrac12m,\tfrac12m+\tfrac12Mi^\lambda,\tfrac12Mi^\lambda$ haben wie $R_1^\lambda$ zum Punkt $0$, bei geradem $\lambda$
\[
\begin{aligned}
R_2^\lambda&=-g''-i^\lambda G''=-\tfrac12(b)-\tfrac12(\mathfrak B),\\
R_3^\lambda&=-i^\lambda G'+i^\lambda g_{\prime}=i^\lambda(\sigma)-\tfrac12i^\lambda(\rho\sigma)-\tfrac12i^\lambda(\rho\sigma ab\mathfrak A\mathfrak B)\\
R_4^\lambda&=+i^\lambda g_{\prime\prime}+G_{\prime\prime}=\tfrac12i^\lambda(b)+\tfrac12i^\lambda(\mathfrak B)
\end{aligned}
\]
und bei ungeradem $\lambda$
\[
\begin{aligned}
R_2^\lambda&=-g''-i^{1-\lambda}H''=-\tfrac12(b)-\tfrac12(\mathfrak B),\\
R_3^\lambda&=-i^{1-\lambda}H'-i^{\lambda-1}h_{\prime}=0,\\
R_4^\lambda&=-i^{\lambda-1}h_{\prime\prime}+G_{\prime\prime}=\tfrac12i^{\lambda+1}(a)+\tfrac12i^{\lambda+1}(\mathfrak A)
\end{aligned}
\]

[V.] Art. $[14.]$ Die Auswertung der Summen der nach Vorschrift III gebildeten $\varepsilon$ ergibt sich aus der Gleichung
\[
\begin{aligned}
&\mathrm{III},\quad\Sigma\varepsilon\text{ von allen }-4\Sigma\varepsilon\text{ von denen, für die }y\text{ ganz, }[x]\text{ gerade}\\
&=\mathrm{III},\quad\Sigma[-\mathrm{Int.}(p-m\omega)+\mathrm{Int.}(p+m\omega)]
\end{aligned}
\]
die anhand der Definition der $\varepsilon$ leicht zu überprüfen ist. Dabei nimmt $p$ alle Werte an, die die unter Vorschrift III angegebenen Bedingungen erfüllen. Mit Ausnahme der beiden Intensoren, die dem kleinsten $(\xi^{*})$ und dem größten zulässigen Wert $(\xi^{**})$ von $\xi$ entsprechen und jeweils gleich
\[
\begin{gathered}
-\mathrm{Int.}(p^{*}-m\omega)\text{ und }+\mathrm{Int.}(p^{**}+m\omega)\text{ oder}\\
-\mathrm{Int.}(\omega m)\text{ und }+\mathrm{Int.}(\tfrac12-\omega)m
\end{gathered}
\]
sind, lassen sich diese Intensoren stets paarweise als $+\mathrm{Int.}(p'+m\omega)$ und $-\mathrm{Int.}(p''-m\omega)$ so zusammenstellen, dass zwischen $\xi'$ und $\xi''$, die den Größen $p'$ und $p''$ entsprechen, kein Wert von $\xi$ liegt, der den reellen oder imaginären Teil von $p$ zu einer ganzen Zahl macht. Die beiden Intensoren heben sich daher immer gegenseitig auf.

\SourceLeaf{394}
[V.] Art. $[15.]$ Es gilt
\[
\begin{aligned}
\mathrm{II},\quad&+\Sigma\varepsilon\mathrm{Int.}\,ip\text{ wobei }Y\text{ ganz }[X]\text{ gerade},\\
&-\Sigma\varepsilon\mathrm{Int.}\,ip\text{ wobei }X\text{ ganz }[Y]\text{ ungerade}\\
={}&\Sigma[-\mathrm{Int.}\,i(p-m\omega)+\mathrm{Int.}\,i(p+m\omega)]\\
&\text{für diejenigen }p,\text{ für die }x\text{ oder }y\text{ ganz, }[X]\text{ gerade}\\
&Y\text{ ungerade, }0<\xi<\tfrac12,\quad\eta=\omega\\
&+\mathrm{Int.}\,i(p^{*}-m\omega)\text{ wenn }[X^{*}]\text{ gerade }[Y^{*}]\text{ ungerade}\\
&-\mathrm{Int.}\,i(p^{**}+m\omega)\text{ wenn }[X^{**}]\text{ gerade }[Y^{**}]\text{ ungerade}
\end{aligned}
\]
wie man leicht erkennt, wenn man auf der zweiten Seite der Gleichung nach dem Verfahren der vorhergehenden Anmerkung über jedes Intervall von $\xi_{\prime}$ bis $\xi_{\prime\prime}$ summiert, das so klein ist, dass zwischen $\xi_{\prime}$ und $\xi_{\prime\prime}$ kein $\xi$ liegt, das im zugehörigen $P$ den reellen oder imaginären Teil zu einer ganzen Zahl macht. Die Anwendung der nach Vorschrift III gebildeten $\varepsilon$ bringt die zweite Seite dieser Gleichung in die in Artikel $15$ angegebene Form.

[V.] Art. $[15.]$ Die Umformung der Summen der nach Vorschrift IV gebildeten $\lambda'\varepsilon$ in die Summen der $\varepsilon$ aus V ergibt sich durch dieselben Überlegungen wie in der letzten Anmerkung, wenn man außerdem die Gleichung
\[
\begin{aligned}
&\mathrm V,\quad\Sigma\varepsilon\text{ von allen, }-4\Sigma\varepsilon\text{ wobei }x\text{ ganz }[y]\text{ gerade},\\
&=+\mathrm{Int.}(\tfrac12im+m\omega)-\mathrm{Int.}(\tfrac12im+\tfrac12-\omega m)
\end{aligned}
\]
verwendet, die der zuvor ermittelten Auswertung der Summe der $\varepsilon$ in Vorschrift III entspricht.

[V.] Art. $[17.]$ Bestimmt man die Hilfsgrößen $U$, $T$, $L$, $V$ durch die Gleichungen
\[
\begin{aligned}
U={}&1-2(B)-(AB)+(\alpha a)+(\beta b)+(\alpha b)-(\beta a)\\
&+(\alpha\beta ab)-(\alpha\beta abAB)\quad\text{oder}\\
U={}&(1+(A))(1-(B))(1-(a\beta)+(b\beta)+(\alpha\beta))
\end{aligned}
\]
\[
\begin{gathered}
\text{weil }(\alpha a)+(\beta b)=(A)+(A\alpha\beta ab),\\
(\alpha b)-(\beta a)=(B)-(B\alpha\beta ab)\text{ ist},
\end{gathered}
\]
\[
\begin{aligned}
T&=-2+(a)+(b)-(\beta)-2(B)-(\alpha ab)\quad\text{oder}\\
T&=-2+(a)+(b)-(\beta)-2(B)-(aA)+(bB)-(\beta AB)\\
L&=(1+(A))(1+(B))(1+(\beta))-(1-(A))(1-(B))(1-(\beta))\\
V&=(1+(A))(1-(B))(-2(a)+2(\beta)+(ab)+(\alpha\beta))\quad\text{oder}\\
V&=(1+(A))(1-(B))(-(a)+(b)+(\beta)+(ab)-(\beta ab)+(\alpha\beta))
\end{aligned}
\]
\[
\text{weil }(a)+(b)=(\beta)+(\beta ab)\text{ wenn }A\text{ positiv }B\text{ negativ}
\]
und bezeichnet mit $W'$, $S'$, $Q'$ die Größen, in die $W$, $S$, $Q$ im Ausdruck für $\mathrm{Dec.}\frac mM$ in Artikel $14$ übergehen, wenn man darin $m$ mit $M$, also $\alpha+\beta i$ mit $\alpha-\beta i$, vertauscht, so gilt

\SourceLeaf{395}
\[
\begin{aligned}
2W'&=-5(B)-(AB)\quad\text{wenn }\frac{M-1}{2}\text{ gerade}\\
2W'&=-(B)-(AB)\quad\text{wenn }\frac{M-1}{2}\text{ ungerade}\\
2S'&=-(\alpha\beta abAB)-(\beta)-(B)+(b)\\
2Q'+2S'+2W'&=2T+U\quad\text{wenn }\frac{M-1}{2}\text{ gerade}\\
2Q'+2S'+2W'&=-4+4(a)+U\quad\text{wenn }\frac{M-1}{2}\text{ ungerade}\\
8\psi&=8(q+r+s+w)-(2Q'+2S'+2W')
\end{aligned}
\]

Ersetzt man hier $8(q+r+s+w)$ durch den in Artikel $15$ angegebenen Wert, bringt ihn jedoch in die Form
\[
\begin{array}{lll}
2T+4L+V&\text{wenn }\dfrac{M-1}{2}\text{ gerade}&\dfrac{m-1}{2}\text{ gerade}\\[4pt]
2T-16+16(A)+V&\text{wenn }\dfrac{M-1}{2}\text{ gerade}&\dfrac{m-1}{2}\text{ ungerade}\\[4pt]
-4+4(a)+V&\text{wenn }\dfrac{M-1}{2}\text{ ungerade}&\dfrac{m-1}{2}\text{ gerade}\\[4pt]
-20+4(a)+V&\text{wenn }\dfrac{M-1}{2}\text{ ungerade}&\dfrac{m-1}{2}\text{ ungerade}
\end{array}
\]
und beachtet, dass
\[
V-U=-\tfrac12(1+(a))(1+(A))(1+(\alpha))(1-(b))(1-(B))(1-(\beta))
\]
gilt, so erhält man für $\psi$ die in Artikel $17$ angegebene Bestimmung.

[VI.] Art. $3$. Das unvollständige Zitat kann auf Artikel $4$ des Bruchstücks III bezogen werden.
\begin{flushright}\textsc{Schering.}\end{flushright}
\BookRule{.3\textwidth}
